\section{Introduction}
Positron emission tomography (PET) maps radiotracer activity through coincident photon detection. Reducing injected activity or acquisition duration lowers the expected number of detected events and increases stochastic variation. A Poisson model remains a useful first approximation after common corrections and normalization, although reconstructed clinical images also contain spatial correlations and modelling errors \citep{shepp1982maximum,qi2006iterative}. Low-dose PET preprocessing must therefore suppress count-dependent fluctuations without biasing uptake values or erasing lesion boundaries.

Conventional Gaussian filtering does not account for signal-dependent statistics. Non-local means (NLM) exploits repeated patches, but patch similarity becomes unreliable under severe noise \citep{buades2005nonlocal}. Total variation (TV) preserves discontinuities but can introduce staircase artefacts \citep{rudin1992nonlinear}. This study addresses these limitations through a variance-stabilized objective that adapts first- and second-order regularization to local image structure.

The main contributions are: (i) a variance-stabilized formulation for low-count PET-like data; (ii) an edge-adaptive combination of first- and second-order penalties; (iii) a non-negative projected gradient solver with Barzilai--Borwein spectral steps and Armijo backtracking; and (iv) a reproducible evaluation including baseline comparison, convergence analysis, ablation, and parameter sensitivity. All reported numerical values are generated from the saved experiment outputs.
